% dcrgraph-doc.tex -- the manual of the dcrgraph package.
% Compile with LuaLaTeX (the import examples need it), twice, in this
% folder (manual/). The examples are in examples/: each is shown as code
% and executed, so the code shown is exactly what made the picture.
%
% Structure (decided with the author):
%   1 Introduction     what the package is, the two ways in, requirements
%   2 Drawing with TikZ  by topic, each with an example; every topic that
%                      applies to imported graphs too is explained here only
%   3 Importing XML    first import, options, relation lists, customising
%                      the imported graph (\adjust...), export, what is read
%   4 Reference        all commands, keys and settings
%   Appendix           DCR portal compatibility
% Every example is shown as code and as its result, from the same source.
\documentclass[a4paper,11pt]{article}
\usepackage[margin=2.2cm]{geometry}
\usepackage[T1]{fontenc}
\usepackage{lmodern}
\usepackage[scaled=0.92]{helvet}
% dcrgraph.sty is in the folder above
\makeatletter\def\input@path{{../}{./}}\makeatother
\usepackage{dcrgraph}
\usepackage{booktabs,tabularx,array}
\usepackage{tcolorbox}
\usepackage{listings}
\usepackage[hidelinks]{hyperref}
\usepackage{parskip}   % space between paragraphs, no indent
\sloppy                % long \texttt names: break lines loosely

\lstset{basicstyle=\ttfamily\small, columns=fullflexible, keepspaces,
  upquote=true}
% \example{name}: examples/name.tex as code, and below it what it produces
% (not tcblisting: running dcrgraph code through its listing mechanism
% fails, the environment's commands stay undefined)
\newcommand\example[1]{%
  \begin{tcolorbox}[colback=white, colframe=black!25, boxrule=0.4pt,
      arc=1mm, left=2mm, right=2mm, halign lower=center]
    \lstinputlisting{examples/#1.tex}
    \tcblower
    \input{examples/#1}
  \end{tcolorbox}}
\newcommand\cs[1]{\texttt{\textbackslash#1}}
\newcommand\key[1]{\texttt{#1}}
\newcommand\meta[1]{\ensuremath{\langle}\textit{#1}\ensuremath{\rangle}}
\newcommand\marg[1]{\texttt{\{}\meta{#1}\texttt{\}}}
\newcommand\oarg[1]{\texttt{[}\meta{#1}\texttt{]}}
\newcommand\draft[1]{\par\noindent\textit{[Draft: #1]}\par}
\newcolumntype{L}{>{\raggedright\arraybackslash}X}

\title{The \textsf{dcrgraph} package\\[4pt]
  \large Drawing DCR graphs in \LaTeX}
\author{tilzuck}
\date{Version 0.2 -- draft, subject to change}

\begin{document}
\maketitle
\tableofcontents

% =====================================================================
\section{Introduction}

\textsf{dcrgraph} draws DCR (Dynamic Condition Response) graphs in
\LaTeX, in the look of the DCR portal of DCR Solutions or in the classic
notation of the literature. There are two ways to create a graph:
\begin{itemize}
  \item \textbf{Drawing with Ti\emph{k}Z} (Section~\ref{sec:tikz}): you write
    the events and relations yourself, in a \texttt{dcrgraph} environment.
  \item \textbf{Importing XML} (Section~\ref{sec:import}): you model the graph
    in the DCR portal (or DCR-js), export it as XML and include the file;
    the graph looks as in the portal, and can be customised in \LaTeX.
\end{itemize}

\paragraph{Requirements.}
\begin{itemize}
  \item Only the file \texttt{dcrgraph.sty} and \cs{usepackage\{dcrgraph\}}.
    Everything else it needs is part of every \TeX{} distribution.
  \item Graphs drawn with Ti\emph{k}Z work with pdf\LaTeX{} and Lua\LaTeX.
  \item Importing needs \textbf{Lua\LaTeX}. On Overleaf: upload
    \texttt{dcrgraph.sty} and choose Menu $\to$ Compiler $\to$ LuaLaTeX.
  \item Graphs with several relations between the same two events need two
    \LaTeX{} runs; latexmk and Overleaf do this automatically.
\end{itemize}

% =====================================================================
\section{\texorpdfstring{Drawing with Ti\emph{k}Z}{Drawing with TikZ}}\label{sec:tikz}

\subsection{A first graph}
As a running example we use the evaluation of applications by a board:
an evaluation round is started; applications can be received until the
deadline has passed; before the board meets, conflicts of interest are
assessed (a subprocess); after the meeting the report is updated and then
approved.

Everything happens in a \texttt{dcrgraph} environment, which is a
Ti\emph{k}Z picture. Events come first, then collections (here the
subprocess, around its events), then relations:
\example{opening}
Events are written \cs{activity}\oarg{options}\marg{label}\marg{id}; the
\meta{id} is used to refer to the event. Relations are written
\meta{relation}\oarg{options}\marg{from}\marg{to}, always from the source
to the target.

\subsection{Events and marking}
The marking (pending, executed, excluded) and the type of an event are
keys of \cs{activity}; \key{role=} and \key{subtitle=} add text.
\example{events}
Events are placed with Ti\emph{k}Z: relative to others
(\key{right=1.5cm of a}, \key{below=2cm of a}) or at coordinates
(\key{at=\{(4,0)\}}). Their size is \key{minimum width=}, \key{minimum height=}
(default 2.6\,cm $\times$ 2.8\,cm).

\subsection{Relations}
All relations at a glance (portal notation; the marker sits at the target
for conditions, pre-conditions and milestones, at the source for all
others):
\example{relations-all}

\paragraph{Placing the lines.}\mbox{}\par
A relation is a straight line between the two boxes. Several relations
between the same two events are spread out side by side. A line can start
and end at a given point of a box (\key{a.south}), run through points
(\key{via=}, which keeps right angles), or take any Ti\emph{k}Z option.
A relation from an event to itself is a marker on its top edge.
\example{routing}

\subsection{Guards and time}\label{sec:values}
A relation can have a guard (\key{guard=\{}\meta{expression}\key{\}}) and a
time (\key{delay=}\meta{time}: for a response a deadline, for the others a
delay). In the portal notation they are shown as markers: a shield for a
guard, a clock for a time. Their values can be written as text, for the
whole graph or for single relations:
\example{values}
\begin{itemize}
  \item For the graph: \key{dcr/guards=}, \key{dcr/delays=},
    \key{dcr/deadlines=} with \key{explicit} (the value as text), \key{none}
    (nothing at all, not even the marker), or the default (\key{shield},
    \key{marker}); \key{dcr/explicit} for all three.
  \item For one relation: \key{explicit=guard}, \key{explicit=\{delay,guard\}},
    \key{explicit} (all).
\end{itemize}
In the classic notation there are no shields or clocks: there the values
are written as text by default, and \key{none} leaves them out.

\subsection{Collections}
Nestings, subprocesses (also multi-instance), forms and subgraphs group
events. They are written after their events and placed around them;
relations can start and end at a collection like at an event.
\example{collections}

\subsection{Portal and classic notation}
The relations can be drawn with the markers of the DCR portal (the
default) or with the arrows of the literature:
\example{notation}
The classic notation has no symbols for two relations of the DCR portal;
they are drawn as what they mean: a pre-condition as a condition plus a
milestone (a time on the condition only), a logical include with the guard
$g$ as an include with the guard $g$ plus an exclude with the guard
\texttt{not(}$g$\texttt{)} (format: \cs{dcrnotguard}). The notation of a whole
document: \cs{usepackage[classic]\{dcrgraph\}}; of one graph:
\key{dcr/notation=classic}.

\subsection{Customisation}
\paragraph{Colours.} \key{bgcolor}, \key{bordercolor} and \key{textcolor} take
any \textsf{xcolor} colour; own colours are defined with \cs{definecolor}.
\example{colours}
The defaults of a whole document: \cs{dcrset\{bgcolor=\dots, nesting
bgcolor=\dots\}}, \cs{colorlet\{dcr border\}\{\dots\}}; the colours of the
relations: \cs{colorlet\{dcr condition\}\{\dots\}} and so on
(Section~\ref{sec:colours}).

\paragraph{Sizes and fonts.} Markers, the spacing of parallel lines, line
widths and the font are settings of the graph or the document
(Section~\ref{sec:settings}); boxes have any size.
\example{sizes}

\subsection{Symbols in running text}
For papers: $A \dcrcondition B$, $A \dcrresponse B$, $A \dcrinclude B$,
$A \dcrexclude B$, $A \dcrmilestone B$, written
\verb|$A \dcrcondition B$| and so on (Section~\ref{sec:symbols}).

% =====================================================================
\section{Importing XML}\label{sec:import}

\subsection{A first import}
The running example of Section~\ref{sec:tikz}, the evaluation of
applications, modelled in DCR-js and exported as XML. Put the file next to
your document and include it:
\example{import-default}
The graph looks as in the editor: positions and sizes, colours, roles,
marking, all events, collections and relations with their markers.
\cs{includedcrgraph} draws an ordinary \texttt{dcrgraph} picture, so
everything in Section~\ref{sec:tikz} applies to it too: the notation, the
display of guards and times, the settings.

\subsection{Options}
The options of \cs{includedcrgraph} (Section~\ref{sec:include}) choose the
notation (\key{notation=classic}), how guards and times are shown, and a
few things of the import itself. Relations are drawn as straight lines by
default; files from DCR-js also store the routes of the lines, which
\key{waypoints} uses:
\example{import-waypoints}
Guards, delays and deadlines can also be
shown for single relations, named by their ids in the file as
``\meta{source} \meta{type} \meta{target}'' (or ``\meta{source}
\meta{target}'' if there is only one relation between the two; \key{value}
is the update relation):
\example{import-options}
A name that matches no relation, or a relation without such a value, gives a
warning.

\subsection{Customising the imported graph}\label{sec:adjust}
The imported graph can be changed in \LaTeX{} without touching the file:
put the changes into a \texttt{dcrgraphfile} environment. Events and
relations are named by their ids in the file (DCR-js generates ids like
\texttt{Event\_191im4m}; the DCR portal uses the ids given in the editor).
\example{adjust}
\begin{itemize}
  \item \cs{adjustrelation}\marg{relation}\marg{options}: \key{from=} and
    \key{to=} set the point of the box where the line starts and ends
    (\key{north}, \key{south}, \key{east}, \key{west}, \key{south east},
    \dots, or an angle like \key{-60}); any other option of a relation
    works too: \key{bend left=20}, \key{via=\{(5,-11)\}}, \key{explicit=guard}.
  \item \cs{adjustactivity}\marg{id}\marg{options}: \key{shift=\{(x,y)\}},
    \key{xshift=}, \key{yshift=} move a box (a collection with everything
    inside it, and the collections around it grow); any other option of an
    activity works too (colours, \key{subtitle=}, \dots).
\end{itemize}
When the graph is changed in the portal and exported again, the changes
still apply; one that names something no longer in the file gives a warning.

\subsection{Taking over the code}
With \key{export=}\meta{file.tex}, \cs{includedcrgraph} also writes the
\texttt{dcrgraph} code of the graph into a file (and into the \texttt{.log}).
Include it with \cs{input} instead of the XML, and edit it freely; it
also works with pdf\LaTeX. On Overleaf the file is under ``Logs and output
files''; copy it into your project.

\subsection{What is read}
From DCR portal files: events with labels, roles, types, marking and
colours; nestings, subprocesses, forms and subgraphs (with their own
events); all relations with their guards and times. The positions are
taken from the file; the portal's own routing of lines is not (relations
are straight lines, spread out where they would overlap). Things the
package cannot draw are skipped with a warning. Details:
Appendix~\ref{sec:compat}.

% =====================================================================
\section{Reference}

\subsection{The environment}
\begin{tabularx}{\linewidth}{@{}lL@{}}
\toprule
\cs{begin\{dcrgraph\}}\oarg{options} & A Ti\emph{k}Z picture; options:
  any Ti\emph{k}Z option, and the settings of Section~\ref{sec:settings} with
  the prefix \key{dcr/} (e.g. \key{dcr/notation=classic}).\\
\bottomrule
\end{tabularx}

\subsection{Events}
\cs{activity}\oarg{options}\marg{label}\marg{id}

\medskip
\begin{tabularx}{\linewidth}{@{}lL@{}}
\toprule
Key & Meaning\\
\midrule
\key{role=}\meta{text} & role in the header; several: \key{role=\{A, B\}}\\
\key{subtitle=}\meta{text} & grey italic text below the label, cut off with \dots\\
\key{pending} & blue ``!'' top left\\
\key{executed} & green tick top left\\
\key{excluded}, \key{included} & dashed outline / normal (default)\\
\key{data} & data event: page icon in the top right corner\\
\key{computation} & computation event: blue ``='' top right\\
\key{dmn} & DMN event (computation by decision table): blue table grid\\
\key{global} & global event: two squares bottom left\\
\key{access} & access control: grey ``\#'' top right\\
\key{form} & event with a form: ``f'' bottom right\\
\key{subgraph} & subgraph (template) event, target of spawns: dotted, ``$|||$'', ``g''\\
\key{bgcolor=}, \key{bordercolor=}, \key{textcolor=} & colours of this event\\
anything else & passed to the Ti\emph{k}Z node: \key{at=\{(x,y)\}},
  \key{right=2cm of a}, \key{minimum width=3cm}, \dots\\
\bottomrule
\end{tabularx}

The label is wrapped to the width of the box; a box grows taller (evenly
around its centre) when its label needs more room than it has, so long labels
never run out of the box.

\subsection{Collections}
Written after their events; \meta{children} is a list of ids.

\medskip
\begin{tabularx}{\linewidth}{@{}lL@{}}
\toprule
Command & Meaning\\
\midrule
\cs{nesting}\oarg{options}\marg{label}\marg{id}\marg{children} & nesting, ``n''\\
\cs{subprocess}\oarg{options}\marg{label}\marg{id}\marg{children} & subprocess, ``s'';
  with \key{multi}: multi-instance\\
\cs{form}\oarg{options}\marg{label}\marg{id}\marg{children} & form, ``f''\\
\cs{subgraph}\oarg{options}\marg{label}\marg{id}\marg{children} & subgraph with its
  own events: dotted, ``$|||$'', ``g''\\
\midrule
Keys & \key{role=}, \key{subtitle=}, \key{excluded}, \key{bgcolor=},
  \key{bordercolor=}, \key{textcolor=}, \key{multi};
  \key{area=\{(x1,y1)(x2,y2)\}}: exact box instead of fitting the children\\
\bottomrule
\end{tabularx}

\subsection{Relations}
\meta{relation}\oarg{options}\marg{from}\marg{to}

\medskip
\begin{tabularx}{\linewidth}{@{}llL@{}}
\toprule
Command & Colour & Marker (portal notation)\\
\midrule
\cs{condition} & orange & open circle at the target\\
\cs{precondition} & orange & filled dot at the target\\
\cs{milestone} & purple & circled ``!'' at the target\\
\cs{response} & blue & ``!'' square at the source\\
\cs{include} & green & ``+'' square at the source\\
\cs{exclude} & red & ``$-$'' square at the source\\
\cs{logicalinclude} & green & ``$\pm$'' square at the source\\
\cs{noresponse} & brown & ``$\times$'' square at the source\\
\cs{update} & grey & ``='' square at the source (the ``value'' relation)\\
\cs{spawn} & black & ``$\ast$'' square at the source (to a subgraph)\\
\bottomrule
\end{tabularx}

\medskip
\begin{tabularx}{\linewidth}{@{}lL@{}}
\toprule
Option & Meaning\\
\midrule
\key{delay=}\meta{time} & timed relation: clock marker; for a response a deadline\\
\key{guard}, \key{guard=\{}\meta{expr}\key{\}} & guard: shield, or the expression\\
\key{explicit=\{guard,delay,deadline\}} & write these values of this relation as
  text; \key{explicit} alone: all\\
\key{via=\{(x1,y1),(x2,y2)\}} & route the line through these points\\
\key{selfpos=}\meta{length} & for a relation from an event to itself: where its
  marker sits on the top edge (its centre, from the left edge); without it,
  the markers are lined up from the left corner (\cs{includedcrgraph} sets
  it, keeping clear of other relations)\\
\key{bundle=}\meta{node}\key{:}\meta{name} & spread this relation together with
  the others of the same bundle, as if they were parallel; for relations
  with a common end \meta{node} whose lines or markers would lie on top of
  each other (\cs{includedcrgraph} sets this automatically)\\
anything else & Ti\emph{k}Z: \key{bend left=20}, \key{out=90}, \dots\\
\bottomrule
\end{tabularx}

\medskip\noindent
A relation from an event to itself is drawn as a marker above the event
(portal) or a loop (classic). Several relations between the same two events
are spread out automatically, unless they have Ti\emph{k}Z options.

\subsection{Settings}\label{sec:settings}
With \cs{dcrset}\marg{settings} for the document, or as \key{dcr/}\meta{setting}
in the options of one graph.

\medskip
\begin{tabularx}{\linewidth}{@{}lL@{}}
\toprule
Setting & Meaning (default)\\
\midrule
\key{notation=portal|classic} & markers of the DCR portal, or arrows (portal)\\
\key{guards=shield|explicit|both|none} & how guards are shown\\
\key{delays=marker|explicit|none} & times of conditions etc.\\
\key{deadlines=marker|explicit|none} & times of responses\\
\key{explicit} & all three explicit\\
\multicolumn{2}{@{}p{\linewidth}@{}}{\small Defaults: portal notation shows only
  the markers (shield, clock); classic notation has no shields or clocks and
  shows the values as text. \key{explicit} adds the text (portal), \key{none}
  shows nothing at all (both notations).}\\
\key{font=} & font of the graph (\cs{sffamily})\\
\key{subtitle font=} & font of subtitles\\
\key{bgcolor=}, \key{nesting bgcolor=} & default backgrounds\\
\key{header height=} & height of the header (0.45\,cm)\\
\key{border width=} & outline of boxes (1.5\,pt)\\
\key{relation width=} & relation lines (1.3\,pt)\\
\key{marker size=} & markers (10.3\,pt); separately:
  \key{condition marker size=}, \key{response marker size=}\\
\key{marker gap=} & gap between marker and box (0\,pt)\\
\key{parallel distance=} & distance of parallel relations (13.5\,pt)\\
\key{classic marker size=} & symbols in the classic notation\\
\key{data size=}, \key{dmn size=}, & sizes of the icons\\
\key{guard size=}, \key{tick size=} & \\
\bottomrule
\end{tabularx}

\medskip\noindent
Package options: \cs{usepackage[classic]\{dcrgraph\}} (or \key{portal}).

\subsection{Colours}\label{sec:colours}
Change with \cs{definecolor} or \cs{colorlet}, e.g.
\cs{colorlet\{dcr condition\}\{orange\}}.

\medskip
\begin{tabularx}{\linewidth}{@{}>{\raggedright\arraybackslash}p{0.6\linewidth}L@{}}
\toprule
Colour & Used for\\
\midrule
\key{dcr activity}, \key{dcr nesting} & backgrounds (\#F9F7ED)\\
\key{dcr border}, \key{dcr text} & outlines (\#CCCCCC), labels (black)\\
\key{dcr condition}, \key{dcr precondition}, \key{dcr response},
\key{dcr include}, \key{dcr exclude}, \key{dcr logical include},
\key{dcr milestone}, \key{dcr noresponse}, \key{dcr update},
\key{dcr spawn} & relations\\
\key{dcr executed}, \key{dcr delay}, \key{dcr subtitle}, \key{dcr icon},
\key{dcr dmn} & markers, labels and icons\\
\key{dcr multi}, \key{dcr multi border}, \key{dcr subgraph},
\key{dcr subgraph collection} & multi-instance subprocesses, subgraphs\\
\bottomrule
\end{tabularx}

\subsection{Importing: \texorpdfstring{\cs{includedcrgraph}}{includedcrgraph}}\label{sec:include}
\cs{includedcrgraph}\oarg{options}\marg{file} \quad (Lua\LaTeX)

\medskip
\begin{tabularx}{\linewidth}{@{}lL@{}}
\toprule
Option & Meaning\\
\midrule
\key{waypoints} & draw relations along the routes stored in the file
  (default: straight lines)\\
\key{notation=classic} & classic notation\\
\key{guards=}, \key{delays=}, \key{deadlines=} & as the settings above, or a list of
  relations: \key{guards=\{A4 response A6, A6 A7\}}; each entry
  ``\meta{source} \meta{type} \meta{target}'' or ``\meta{source} \meta{target}''
  (ids from the file; \key{value} = \cs{update})\\
\key{explicit} & all values as text\\
\key{marking=false} & without executed / pending / excluded\\
\key{scale=}\meta{factor} & scale the positions (not the boxes)\\
\key{export=}\meta{file.tex} & also write the generated code to a file, and
  into the \texttt{.log}\\
\key{format=dcrjs} & input format (the only one so far)\\
\bottomrule
\end{tabularx}

\subsection{Adjusting an imported graph}
\begin{tabularx}{\linewidth}{@{}>{\raggedright\arraybackslash}p{0.47\linewidth}L@{}}
\toprule
Command & Meaning\\
\midrule
\cs{begin\{dcrgraphfile\}}\oarg{options}\marg{file} & an imported graph with
  adjustments; options as for \cs{includedcrgraph}; drawn at
  \cs{end\{dcrgraphfile\}} (Lua\LaTeX)\\
\cs{adjustactivity}\marg{id}\marg{options} & \key{shift=\{(x,y)\}},
  \key{xshift=}, \key{yshift=} (in cm, or with a unit): move the box, a
  collection with everything inside it; other options: as for \cs{activity}\\
\cs{adjustrelation}\marg{relation}\marg{options} & \meta{relation}:
  ``\meta{source} \meta{type} \meta{target}'' or ``\meta{source} \meta{target}'';
  \key{from=}, \key{to=}: anchor of the box where the line starts/ends;
  other options: as for relations\\
\bottomrule
\end{tabularx}

\subsection{Symbols in text}\label{sec:symbols}
In math mode: \cs{dcrcondition} $\dcrcondition$, \cs{dcrresponse} $\dcrresponse$,
\cs{dcrinclude} $\dcrinclude$, \cs{dcrexclude} $\dcrexclude$,
\cs{dcrmilestone} $\dcrmilestone$, e.g.\ \verb|$A \dcrcondition B$|.
The format of guards on lines: \cs{dcrguardformat}; of the negated guard of a
logical include in the classic notation: \cs{dcrnotguard}.

% =====================================================================
\appendix
\section{DCR portal compatibility}\label{sec:compat}
\paragraph{Read from the file.} Events: label, roles, position, colours
(background, border, text), marking (executed, pending, included), type
(data, computation, DMN, global, form, subgraph). Collections: nestings,
subprocesses, forms, subgraphs with their own graph; their boxes are
computed as in the portal. Relations: condition, pre-condition, milestone,
response, include, exclude, logical include, no-response, value (update),
spawn; their guards and times; the routes of lines (only with the option
\key{waypoints}, for DCR-js files).

\paragraph{Not reproduced.} The portal's routing of lines (it draws some
lines at right angles, the package draws straight lines and spreads out
lines that would overlap); the order of parallel lines between two events;
the value assigned by an update relation; data types, expressions and
everything that only concerns the execution of the graph.

\paragraph{Sizes.} Portal boxes are 115\,$\times$\,125\,px; the package draws
them as 2.6\,$\times$\,2.8\,cm (its default box), the positions scaled
accordingly.

\end{document}
